% !TEX root = ../math600-notes.tex

\section{Introduction to Algebraic Geometry}

We study the solutions of systems of polynomials, working over $\C$, although most everything we say can also be done over arbitrary fields.

\begin{defn}
    \textbf{Affine $n$-space} $\A^n:= \A^n(\C)$ is just $\C^n$; its elements are points. $\A^1$ is the affine line, $\A^2$ is the affine plane.

    Normally, the letter $\A$ is used to denote an arbitrary field, but we're only looking at $\C$, so we may just think of it as $\C^n$.

    Given $f \in \C[x_1,\dots,x_n]$, a point $a=(a_1,\dots,a_n) \in \A^n$ is a \textbf{zero} of $f$ if $f(a)=0$. If $f$ is not constant, the set of zeroes of $f$ is the \textbf{hypersurface} defined by $f$, denoted $V(f)$. If $f$ has degree $1$, then $V(f)$ is a \textbf{hyperplane} in $\A^n$.
\end{defn}
\begin{ex}
    The following are pictured below
    \begin{enumerate}
        \item $V(x^2-y^2) \subset \A^2$
        \item $V(y^2-x^2(x+1)) \subset \A^2$
    \end{enumerate}
\end{ex}

\begin{figure}[H]
  \centering
  \begin{minipage}[b]{0.3\textwidth}
    \includegraphics[width=\textwidth]{ex1.png}
  \end{minipage}
  \hfill
  \begin{minipage}[b]{0.5\textwidth}
    \includegraphics[width=\textwidth]{ex2.png}
  \end{minipage}
\end{figure}

\begin{defn}
    More generally, if $S$ is any set of polynomials in $\C[x_1,\dots,x_n]$, we let
    $$V(S) = \{ P \in \A^n \mid f(P) =0 ,\forall f \in S\}=\bigcap_{f \in S}V(f)$$
    If $S = \{f_1,\dots,f_r\}$, we write $V(f_1,\dots,f_r)$ instead of $V(\{f_1,\dots,f_r\})$. A subset $X \subset \A^n$ is an \textbf{affine algebraic set} if $X = V(S)$ for some $S$.
\end{defn}

\begin{prop}
    \begin{enumerate}
        \item If $I=(S)$ is the ideal of $\C[x_1,\dots,x_n]$ generated by $S$, then $V(S) = V(I)$. So, every algebraic set is equal to $V(I)$ for some ideal $I$
        \item If $\{I_\alpha\}$ is any collection of ideals, then
        $$V(\bigcup_{\alpha}I_\alpha) = \bigcap_{\alpha}V(I_\alpha).$$
        So, the intersection of algebraic sets is an algebraic set.
        \item If $I \subset J$, then $V(J) \subset V(I)$.
        \item $V(fg) = V(f) \cup V(g)$ for all $f,g \in \C[x_1,\dots,x_n]$. More generally, $V(IJ) = V(I) \cup V(J)$. So, any finite union of algebraic sets is an algebraic set.
        \item $V((0)) = \A^n$, $V((1))=\emptyset$
        \item[(5')] $V(X_1-a_1,X_2-a_2,\dots,X-a_n)=\{(a_1,\dots,a_n)\}$
    \end{enumerate}
    Properties (2), (4), and (5) say that the algebraic sets are the closed sets of a topology on $\A^n$. This is the \textbf{Zariski topology}. (5') show that finite subsets of $\A^n$ are algebraic.
\end{prop}

\begin{notebox*}[Exercise]
    Check that $V(IJ)=V(I \cap J)$, even though $IJ \subsetneq I \cap J$ in general.
\end{notebox*}

\begin{ex}
    In $\A^2$, coordinates $x$ and $y$, let $I:= (x),J:=(y)$, $IJ=(xy),\angles{I \cup J}=\angles{x,y}$. $V(x)$ is the $y$-axis, $V(y)$ is the $x$-axis, $V(xy)$ is the $x$ and $y$ axes, and $V(x,y)$ is the origin.
\end{ex}

\section{The Ideal of a Set of Points}

\begin{defn}
    For any subset $X \subset \A^n$, let
    $$I(X) := \{ f \in \C[x_1,\dots,x_n] \mid f(x) = 0 ,\forall x \in X\}.$$
    Then, $I(x)$ is an ideal of $\C[x_1,\dots,x_n]$ called the ideal of $X$.
\end{defn}

\begin{prop}
    \begin{enumerate}
        \item If $X \subset Y$, then $I(Y) \subset I(X)$
        \item $S \subset I(V(S))$ for any set $S$ of polynomials. Similarly, $X \subset V(I(X))$ for any set $X$ of points.
        \item $V(I(V(S))) = V(S)$ for any set $S$ of polynomials, and $I(V(I(X)))=I(X)$ for any set $X$ of points.
    \end{enumerate}
    So, if $V$ is an algebraic set, then $V = V(I(V))$, and if $I$ is the ideal of an algebraic set, then $I=I(V(I))$.
\end{prop}

\begin{ex}
    Not all ideals are ideals of algebraic sets. For example, $J:= (x^2) \subset \C[x]$ has $V(J)=\{0\}$, but $I(\{ 0\}) = (x) \neq J$. If $I = I(X)$ and $f^n \in I$ for some $n >0$, then $f \in I$. This says exactly that $I = \rad (I)$, so $I$ is a radical ideal.
\end{ex}

\begin{prop}
    $I(X)$ is a radical ideal for any set $X \subset \A^n$. (Later, we'll show that radical ideals are exactly the ideals of algebraic sets.)
\end{prop}

\begin{thm}[{Hilbert Basis Theorem For $\C[x_1,\dots,x_n]$, Rephrased}]
    Every algebraic set is the intersection of a finite number of hypersurfaces.
\end{thm}

\begin{notebox*}
    Notice how the algebraic set comprised $x$ and $y$ axes together is "reducible", into the irreducible algebraic sets the $x$-axis and $y$-axis.
\end{notebox*}

\begin{defn}
    We say an algebraic set $V$ is \textbf{reducible} if $V = V_1 \cup V_2$, where $V_1$ and $V_2$ are algebraic sets in $\A^n$, $V_i \subsetneq V$, $i=1,2$. Otherwise, we say $V$ is irreducible.
\end{defn}

\begin{prop}
    An algebraic set $V$ is irreducible $\iff$ $I(V)$ is a prime ideal.
\end{prop}

\begin{proof}
    If $I(V)$ is not prime, then there exist $F_1,F_2$ with $F_1F_2 \in I(V)$, but neither $F_1$ nor $F_2$ are in $I(V)$. Then, $V=(V \cap V(F_1)) \cup (V \cap V(F_2))$. Then, $V \cap V(F_1) \neq V$ and $V \cap V(F_2) \neq V$, so $V$ is reducible.

    Conversely, suppose $V= V_1 \cup V_2$, $V_i \subsetneq V$. Then, $I(V_i) \supsetneq I(V)$ [If $I(V_i)=I(V)$, then $V_i=V(I(V_i))=V(I(V))=V$, a contradiction.] Let $F_i \in I(V_i)$, $F_i \not \in I(V)$. Then $F_1F_2 \in I(V)$, so $I(V)$ is not prime.
\end{proof}

\begin{thm}
    If $I$ is a prime ideal, then $V(I)$ is irreducible.
\end{thm}
This is much harder to show.

\begin{thm}
    Let $V \subset \A^n$ be an algebraic set. Then, there are unique irreducible algebraic sets $V_1,\dots,V_m$ such that $V=V_1 \cup V_2 \cup \cdots \cup V_m$, and $V_i \not \subset V_j$ for all $i \neq j$. The $V_i$ are the \textbf{irreducible components} of $V$.
\end{thm}

\begin{proof}
    Recall that since $\C[x_1,\dots,x_n]$ is Noetherian, any collection of ideals has a maximal element. Rephrased, this implies that any collection of algebraic sets has a minimal element. Let $\mathcal S := \{$Algebraic sets $V \subset \A^n \mid$ $V$ is not the union of finitely many irreducible algebraic sets $\}$. We want to show $\mathcal S = \emptyset$. Otherwise, $\mathcal S$ has a minimal member $V$. Since $V \in \mathcal S$, it's not irreducible. So, $V = V_1\cup V_2$ with $\emptyset \neq V_i \subset V$, $V_i$ algebraic. Now, $V_1, V_2 \not \in \mathcal S$, so $V_i = \cup V_{i_k}$ for $i=1,2$ and $V_{i_k}$ irreducibles. Then, $V= \cup_{i,k}V_{i_k}$, a contradiction.

    Hence, any algebraic set $V= V_1 \cup \cdots \cup V_m$ for $V_i$ irreducible. To achieve the second condition, if $V_i \subset V_j$ for some $i \neq j$, then throw away $V_i$. This proves existence. Now, suppose $V=V_1\cup \cdots \cup V_m=W_1\cup \cdots \cup W_s$. Then,
    $$V_i = V_i \cap (\bigcup_{j=1}^s W_j)=\bigcup _{j=1}^s(V_i \cap W_j)$$
    and $V_i$ is irreducible, so $V_i \subset W_{j(i)}$ for some $j(i)$. Similarly, $W_{j(i)}\subset V_k$ for some $k$. Then, $V_i \subset W_{j(i)}\subset V_k \implies k=i$, so $V_i = W_{j(i)}$. Likewise, each $W_j$ is equal to some $V_{i(j)}$.
\end{proof}

\begin{defn}
    Let $R$ be a commutative ring. We define the \textbf{Krull Dimension} of $R$ to be the maximum $m$ such that there exists a chain of prime ideals $P_0 \subsetneq P_1 \subsetneq \cdots \subsetneq P_m \subset R$.
\end{defn}

\begin{notebox*}[Thought Exercise]
    What does the previous theorem say in terms of ideals in $\C[x_1,\dots,x_n]$? More generally, for any commutative Noetherian ring?
\end{notebox*}

\begin{notebox*}[Answer (Lasker-Noether)]
    Primary decomposition theorem for ideals in Noetherian Rings, which generalizes the fundamental theorem of arithmetic.
\end{notebox*}

\section{Algebraic Subsets of the Line and Plane}

For the affine line $\A^1(\C) = \C$, if $V = V(I)$ is an algebraic subset, then $I = (f(x))$, so $V$ is $\emptyset$, $\A^1$, or a finite number of points, which the roots corresponding to the roots of $f$. If $f(x) = \prod_{i=1}^n (x-a_i)^{k_i}$, then $\rad (I) = (\prod_{i=1}^n (x-a_i))$. The irreducible algebraic sets are $\emptyset, \A^1$, and single points.

\begin{prop}
    Let $f,g$ be polynomials in $\C[x,y]$ with no common factors. Then, $V(f,g)=V(f) \cap V(g)$ is a finite number of points.
\end{prop}

\begin{proof}
    $f$ and $g$ are coprime in $\C[x][y]$. By Gauss's lemma, $f$ and $g$ are coprime in $\C(x)[y]$, which is a PID because $\C(x)$ is a field. Hence, $(f,g)=(1)$ in $\C(x)[y]$, so there exists $r,s\in \C(x)[y]$ with $rf+sg=1$. Now, there is a polynomial $d(x) \in \C[x]$ so that $dr=a,ds=b$ both lie in $\C[x,y]$ (eliminating denominators), hence $af+bg=d \in \C[x,y]$. If $(a_0,b_0) \in V(f,g)$, then $d(a_0)=0$. But, $d$ only has a finite number of zeroes, hence there is a finite number of $x$-coordinates which appear in $V(f,g)$. The same reasoning can be applied to the $y$-coordinates, proving that $V(f,g)$ is finite.
\end{proof}

\begin{lem}
    If $F \in \C[x,y]$ is any non-constant polynomial, then $V(F)$ is infinite.
\end{lem}

\begin{proof}
    On homework
\end{proof}

\begin{cor}
    If $f \in \C[x,y]$ is irreducible, then $I(V(f))=(f)$ and $V(f)$ is irreducible.
\end{cor}

\begin{proof}
    If $g \in I(V(f))$, then $V(f,g) \supset V(f)$, which is infinite by the lemma. The previous proposition therefore implies that $f$ and $g$ are not coprime, so as $f$ is irreducible, we have $f \mid g$. That is, $g \in (f)$. The opposite inclusion is immediate, therefore $I(V(f))=(f)$. Finally, since $(f)$ is prime, a previous proposition implies $V(f)$ is irreducible.
\end{proof}

\begin{cor}
    The irreducible algebraic subsets of $\A^2$ are $\emptyset,\A^2,$ points, and irreducible plane curves $V(f)$, where $f \in \C[x,y]$ is irreducible.
\end{cor}

\begin{proof}
    Suppose $V$ is any irreducible algebraic set. If $V$ is finite, or $I(V)=(0)$, then $V$ is of the required type. Otherwise, let $f \in I(V)$ be non-constant, and since $I(V)$ is prime, some irreducible factor of $f$ lies in $I(V)$, so we can assume $f$ is irreducible. Then, we claim $I(V)=(f)$: if $g \in I(V)$, $g \not \in (f)$, then $V \subset V(f,g)$ is finite, because $f,g$ are coprime, a contradiction.
\end{proof}

\begin{cor}
    Let $f \in \C[x,y]$ and $f = \prod _{i=1}^rf_i^{n_i}$ be the decomposition of $f$ into irreducible factors. Then, $V(f) = V(f_1) \cup \cdots \cup V(f_r)$, so this is the decomposition of $V(f)$ into irreducible components, and $I(V(f))=(f_1\cdots f_r)$.
\end{cor}

\begin{proof}
    No $f_i$ divides $f_j$ for $i \neq j$, so there are no inclusion issues among the $V(f_i)$. Moreover,
    $$I(\bigcup _{i}V(f_i))=\bigcap _iI(V(f_i))=\bigcap_i (f_i)=(f_1,\dots,f_r)$$
\end{proof}
\section{Hilbert's Nullstellensatz (Theorem of the Zeroes)}

This gives the relationship between ideals of $\C[x_1,\dots,x_n]$ and algebraic sets.

\begin{notebox*}[Question]
    Given some ideal $I \subset \C[x_1,\dots,x_n]$, when is $V(I) \neq \emptyset$.
\end{notebox*}

If $1 \in I$, then clearly $V(I)$ is empty. What about the converse?

Let $I \subsetneq \C[x_1,\dots,x_n]$ be a proper ideal, and let $M$ be a maximal ideal containing $I$. Then, $V(I) \supset V(M)$. What can we prove about $M$?

\begin{thm}[Weak Nullstellensatz]
    The maximal ideals in $\C[x_1,\dots,x_n]$ are in bijective correspondence with the points of $\A^n=\C^n$. So, each point $(a_1,\dots,a_n)$ corresponds to the ideal
    $$M_a:=(x_1-a_1,\dots,x_n-a_n).$$
\end{thm}

\begin{proof}
    We first check that $M_a$ is indeed maximal. We have an evaluation map at $a=(a_1,\dots,a_n)$: $ev:\C[x_1,\dots,x_n] \to \C,f \mapsto f(a)$. We claim that $\ker (ev)=M_a$. Clearly, $M_a \subset \ker (ev)$. Conversely, assume $f(a)=0$, and expand $f$ into a Taylor expansion around $a$:
    $$f(x)=f(a)+\sum_{i}c_i(x_i-a_i)+\sum_{i,j}c_{i,j}(x_i-a_i)(x_j-a_j)+\cdots$$
    Since $f(a)=0$, we see that $f \in M_a$. Thus, $\C[x_1,\dots,x_n]/M_a \cong \C$ is a field, hence $M_a$ is maximal.

    Now, let $M$ be a maximal ideal of $\C[x_1,\dots,x_n]$. Consider the map
    $$\pi _1:\C[x_1] \hookrightarrow \C[x_1,\dots,x_n] \to C[x_1,\dots,x_n]/M=K,$$
    where $K$ is a field. Then, $\ker \pi _1$ is either $0$ or a maximal ideal of $\C[x_1]$, which we see because $\C[x_1]/\ker \pi_1 \hookrightarrow K$, so it's a domain, and $\C$ is a PID (prime ideals are maximal).

    Now, we claim $\ker \pi_1 \ne 0$. Suppose not: then, $\pi_1 :\C[x] \hookrightarrow K$. By the universal property of the quotient field $\C(x_1)$, $\pi_1$ extends to an injection $\C(x_1) \hookrightarrow K$.
    Thus, $K$ contains a field isomorphic to $\C(x)$. As a vector space over $\C$, $\C[x_1,\dots,x_n]$ has a countable basis, namely $\{x_1^{a_1},\dots x_n^{a_n}\}_{a_i\geq0}$. Therefore $K = \C[x_1,\dots,x_n]/M$ must also have a countable basis $\{v_1,v_2,\dots\}$. Now, if $S \subset K$ is any linearly independent set, then $S$ is also countable, since $S = \bigcup_{n\geq1} (S \cap \angles{v_1,\dots,v_n})$ is a countable union of finite sets. On the other hand, $\C(x)$ contains the elements $\frac{1}{x-\alpha}$, where $\alpha \in \C$, and these form an uncountable linearly independent set, a contradiction. (If $\sum \frac{\lambda_i}{x-a_i}=0$, and $\lambda_i \neq 0$, then the limit as $x \to a_i$ of the LHS is $\infty$). We've therefore shown that $\ker \pi_1 \neq 0$, hence it's a maximal ideal of $\C[x_1]$, which we know must be $(x_1-a_1)$ for some $a_1 \in \C$. It follows that $(x_1-a_1) \in M$, and the same argument can be applied for all $x_i$. So $x_i - a_i \in M$ for $1 \leq i \leq n$, and $a_i \in \C$, but then $M_a \subset M$, so $M=M_a$.
\end{proof}

\begin{cor}
    If $I \subset \C[x_1,\dots,x_n]$ is a proper ideal, then $V(I) \neq \emptyset$.
\end{cor}

\begin{proof}
    $I$ must be contained in some maximal ideal $M$, and $V(I) \supset V(M) \neq \emptyset$
\end{proof}

\begin{cor}
    If $n=1$, we get the fundamental theorem of algebra.
\end{cor}

\begin{thm}[Strong Nullstellensatz]
    Let $I$ be an ideal in $\C[x_1,\dots,x_n]$. Then, $I(V(I))=\rad (I)$. Concretely, this says: Let $I=(f_1,\dots,f_r)$. If $g$ is such that $g$ vanishes wherever $f_1,\dots,f_r$ all vanish, then there is an equation $g^N=a_1f_1+\cdots +a_rf_r$ for some $N \in \N$, $a_1,\dots,a_r \in \C[x_1,\dots,x_n]$.
\end{thm}

\begin{proof}
    Note that $\rad (I) \subset I(V(I))$ is easy (exercise). Suppose $g \in I(V(f_1,\dots,f_r))$, with $f_i \in \C[x_1,\dots,x_n]$. Consider the ideal
    $$J=(f_1,\dots,f_r,x_{n+1}g-1) \subset \C[x_1,\dots,x_n,x_{n+1}].$$
    Then, notice that $V(J)=\emptyset$, because when $f_1,\dots,f_r$ all vanish, $g$ does as well, and then the last component is non-zero. We deduce from weak Nullstellensatz that $1 \in J$, so there exists $a_1,\dots,a_r,b \in \C[x_1,\dots,x_{n+1}]$ with
    $$1 = \sum_{i}a_i(x_1,\dots,x_{n+1})f_i+b(x_1,\dots,x_{n+1})(x_{n+1}g-1)$$
    Now, set $x_{n+1}:= \frac{1}{g}$. This gives
    $$1 = \sum_{i} a_i(x_1,\dots,x_n,\frac{1}{g})f_i,$$
    and multiplying by a large enough $g^N$ to clear the denominators gives
    $$g^N = \sum_ic_i(x_1,\dots,x_n)f_i,\:\:\:c_1,\dots,c_r \in \C[x_1,\dots,x_n]$$
\end{proof}

\begin{cor}[Ideal-Variety Correspondence]
    We'll work in $F[x_1,\dots,x_n]$, where $F$ is any field.
    \begin{enumerate}
        \item The maps $\{$algebraic sets$\} \hookrightarrow \{$ideals$\}$ and $\{$ideals$\} \to \{$algebraic sets$\}$ are inclusion reversing. (i.e. $I_1 \subset I_2 \implies V(I_2)\subset V(I_1)$,etc). Furthermore, for any algebraic set $V$, $V(I(V))=V$, so $I$ is always injective.
        \item If $F=\C$ (or generally $F$ is algebraically closed), then we have inclusion reversing bijections $\{$algebraic sets$\} \longleftrightarrow \{$radical ideals$\}$, containing $\{$irreducible algebraic sets$\} \longleftrightarrow \{$prime ideals$\}$, containing $\{$points$\} \longleftrightarrow \{$maximal ideals$\}$. This is why this is called algebraic geometry, as we have a dictionary between geometric structures and algebraic structures.
    \end{enumerate}
\end{cor}
